\section{Einleitung und Aufgabenstellung}
\label{sec:einleitung}

% VERANTWORTLICH: ...
% WAS HIER HINGEHOERT:
%   - Wozu wird das Modell gebaut? Welche Frage soll es beantworten?
%   - Kurze Nennung der acht Aufgabenpunkte und wo sie im Protokoll stehen.
%   - KEINE Modellgleichungen, die kommen in Abschnitt 2.
% LAENGE: etwa eine Seite.

Die Leistung eines Photovoltaikmoduls haengt neben der Einstrahlung stark von
der Modultemperatur ab. Ziel dieser Arbeit ist ein thermisches Modell, das aus
Globalstrahlung $G$ und Umgebungstemperatur $\Tamb$ die Modultemperatur $\Tm$
und daraus die elektrische Leistung $\Wel$ vorhersagt.

\subsection{Zweck des Modells}
% Diese Frage muss vor allem anderen beantwortet sein. Die Vorlesung
% definiert Modellvalidierung als "erfuellt das Modell seinen Zweck".
% Ohne festgelegten Zweck ist Kapitel 5 nicht schreibbar.

\subsection{Aufbau der Arbeit}
